\documentclass[12pt]{article}
\usepackage[utf8]{inputenc}
\usepackage[T1]{fontenc}
\usepackage[margin=1in]{geometry}
\usepackage{setspace}
\usepackage{xcolor}
\usepackage[colorlinks=true, urlcolor=blue!60!black, linkcolor=blue!60!black]{hyperref}
\setlength{\parindent}{0em}
\setlength{\parskip}{0.6em}
\renewcommand{\baselinestretch}{1.0}
\begin{document}
\begin{center}
{\LARGE \textit{CFPB v. Community Financial Services Ass'n of America, Ltd.}, 601 U.S. 416 (2024)\par}
\vspace{0.5em}
{\large Case Brief --- Two-Page Edition\par}
\end{center}

\section{Memory Jogger}
The CFPB's capped, standing draw on Federal Reserve earnings is a valid appropriation: the Clause requires only a law naming a source of funds and a purpose (\href{https://supreme.justia.com/cases/federal/us/601/22-448/}{Justia}).

\section{Detailed Case Facts}
Dodd-Frank lets the CFPB Director requisition each year from the Federal Reserve's combined earnings the amount ``reasonably necessary'' to run the Bureau, up to an inflation-adjusted cap, outside the annual appropriations process. Payday lenders challenged the Bureau's 2017 Payday Lending Rule.

\section{Procedural History}
The district court upheld the rule; the Fifth Circuit vacated it, holding the ``self-actualizing, perpetual funding mechanism'' unconstitutional.

\section{Judicial Votes}
Reversed 7--2. Thomas wrote for the Court; Kagan (with Sotomayor, Kavanaugh, and Barrett) and Jackson concurred; Alito, joined by Gorsuch, dissented.

\section{Holding}
An appropriation needs only to identify a source of public funds and authorize spending for designated purposes; no annual renewal or fixed sum is required.

\section{Analysis of Opinions}
\textbf{Majority:} founding-era dictionaries, the Army Clause's two-year limit, Parliament's standing and capped grants, and the First Congress's fee-funded Customs Service and Post Office all show flexibility; Congress set the cap. \textbf{Kagan:} two centuries of practice confirm it. \textbf{Jackson:} courts should not invent limits the text lacks. \textbf{Alito:} the CFPB's unprecedented autonomy turns the Clause into ``a minor vestige.''

\section{Examples: Future Applications}
\textbf{Same-Side:} a registry funded by its own filing fees up to a cap. \textbf{Opposite-Side:} a permanent appropriation for the Army violates the two-year limit. \textbf{Fence-Sitter:} an uncapped, perpetual draw for an AI agency, which also raises nondelegation concerns.

\section{Critique}
Critics call the source-and-purpose test toothless; defenders note that Congress can always repeal a standing appropriation and that Alito's ``novelty'' test lacks a textual anchor.

\section{Key Quotations}
\begin{quote}
``Under the Appropriations Clause, an appropriation is simply a law that authorizes expenditures from a specified source of public money for designated purposes.''
\end{quote}
\noindent --- Justice Thomas, for the Court, 601 U.S. at 424.

\end{document}
